\section{Background and Related Work}
\label{sec:related}

\subsection{Inverse Reinforcement Learning}
IRL recovers a reward function from demonstrated behaviour. The foundational
linear-programming formulation~\citep{ng2000algorithms} was followed by the
maximum-entropy formulation~\citep{ziebart2008maxent}, now the standard
probabilistic treatment, and by Bayesian IRL~\citep{ramachandran2007bayesian},
which places a Boltzmann likelihood over trajectories and returns a posterior over
reward weights. IOR adopts the Bayesian formulation precisely because it yields a
posterior rather than a point estimate. IRL is known to be partially identifiable:
many reward functions explain one behaviour set~\citep{partialIdentifiability2024},
which IOR addresses by reporting the full posterior.

\subsection{IRL Applied to Language-Model Alignment}
The Alignment Auditor~\citep{alignmentAuditor2025} applies Bayesian IRL to verify
language-model training objectives, and IR\textsuperscript{3}~\citep{ir3_2026} uses
contrastive IRL to detect reward hacking; related work reconstructs implicit
training goals from behaviour~\citep{insightsInverse2024}. All operate on model
weights or RLHF training signals. IOR differs in requiring only API-level traces,
making it applicable to closed deployments.

\subsection{Agent Misalignment Benchmarks and Red-Teaming}
AgentMisalignment~\citep{agentMisalignment2025} measures the propensity for
misaligned behaviour (sandbagging, power-seeking, oversight avoidance) across
frontier models in agentic settings with tool use, establishing that the problem is
real in deployed agents but not inferring what an agent optimises for. Ulterior
Motives~\citep{ulteriorMotives2026} detects misaligned reasoning in continuous
thought models via probes on latent tokens, requiring model internals. These works
measure \emph{whether} agents misbehave; IOR infers \emph{what} they optimise for.

\subsection{Goal Inference from Behaviour}
Online Bayesian goal inference~\citep{sips2020} infers goals from planning
trajectories in structured environments and shows that inference must survive
non-optimal trajectories; related work infers the goals of communicating
agents~\citep{communicatingAgents2023}. These methods assume structured planning
environments with known transition models. Agent tool-call trajectories lack an
explicit transition model, which is why IOR substitutes the Goal Decomposition
Judge as its feature extractor.

\subsection{Regulatory Context}
The EU AI Act anchors conformity on intended purpose~\citep{euAIAct2024}, and
OWASP~\citep{owaspASI2025}, MITRE ATLAS~\citep{mitreAtlas}, and the NIST AI
RMF~\citep{nistRMF2026} provide the taxonomies onto which IOR maps every finding,
making divergence reports regulatory-legible rather than bespoke.
